\documentclass[12pt, a4paper]{article}

\usepackage{fontspec}
\usepackage{polyglossia}
\setmainlanguage{french}
\usepackage{amsmath, amssymb}
\usepackage{geometry}
\usepackage{booktabs}
\usepackage{listings}
\usepackage{xcolor}
\usepackage{hyperref}
\usepackage{parskip}
\usepackage{algorithm}
\usepackage{algpseudocode}
\usepackage{graphicx}

\geometry{margin=2.5cm}

\definecolor{codebg}{RGB}{245,245,245}
\definecolor{codegreen}{RGB}{0,128,0}
\definecolor{codegray}{RGB}{128,128,128}
\definecolor{codeblue}{RGB}{0,0,200}

\lstset{
  language=Python,
  backgroundcolor=\color{codebg},
  basicstyle=\ttfamily\small,
  keywordstyle=\color{codeblue}\bfseries,
  commentstyle=\color{codegreen},
  stringstyle=\color{orange},
  numberstyle=\tiny\color{codegray},
  numbers=left,
  stepnumber=1,
  frame=single,
  framerule=0.4pt,
  breaklines=true,
  showstringspaces=false,
  tabsize=4,
}

\title{
  \textbf{Savings Algorithm (Clarke-Wright)}\\
  \large Documentation algorithmique --- Travelling Salesman Problem
}
\author{Recherche Opérationnelle --- CESI}
\date{}

\begin{document}

\maketitle
\tableofcontents
\newpage

% =============================================================================
\section{Définition}

L'\textbf{algorithme de Clarke-Wright} (1964) est une heuristique \textbf{constructive}
pour le TSP et le VRP (Vehicle Routing Problem). Il est basé sur la notion
d'\textbf{économie (saving)} : la distance économisée en reliant directement
deux villes $i$ et $j$ plutôt que de passer par le dépôt.

C'est une méthode \textbf{gloutonne} (greedy) : elle construit le tour en fusionnant
itérativement les paires de villes qui maximisent le gain de distance, sans jamais
revenir en arrière. Elle ne garantit \textbf{pas l'optimalité} mais produit des
solutions à 10--15\% de l'optimum en temps $O(n^2 \log n)$.

% =============================================================================
\section{Méthode}

\subsection{Étape 1 --- Configuration initiale}

On désigne une ville comme \textbf{dépôt} (ville $0$). On initialise $n-1$ routes
triviales, chacune reliant le dépôt à une seule ville :

\begin{equation}
  \text{dépôt} \to i \to \text{dépôt} \quad \forall\, i \in \{1, \ldots, n-1\}
\end{equation}

Le coût total de cette configuration initiale est :

\begin{equation}
  C_0 = 2 \sum_{i=1}^{n-1} c_{0i}
\end{equation}

où $c_{0i}$ est la distance entre le dépôt $0$ et la ville $i$.

\subsection{Étape 2 --- Calcul des savings}

Pour chaque paire de villes $(i, j)$ avec $i \neq j$ et $i,j \neq 0$, on calcule
le \textbf{saving} $s(i,j)$ :

\begin{equation}
  s(i, j) = c_{0i} + c_{0j} - c_{ij}
  \label{eq:saving}
\end{equation}

\textbf{Interprétation de \eqref{eq:saving} :} si l'on fusionne les deux routes
$0 \to i \to 0$ et $0 \to j \to 0$ en une seule route $0 \to i \to j \to 0$,
on supprime deux arcs vers le dépôt ($c_{0i}$ et $c_{0j}$) et on ajoute un arc
direct ($c_{ij}$). Le gain net est précisément $s(i,j)$.

Un saving positif ($s(i,j) > 0$) signifie que la fusion réduit la distance totale.
Un saving négatif signifie qu'il vaut mieux conserver les routes séparées.

\subsection{Étape 3 --- Fusion gloutonne}

On trie les savings par ordre \textbf{décroissant} et on parcourt la liste.
Pour chaque paire $(i,j)$, on fusionne si les trois conditions sont réunies :

\begin{enumerate}
  \item $i$ et $j$ appartiennent à des routes \textbf{différentes}
        (pas déjà dans le même tour partiel)
  \item $i$ est une \textbf{extrémité} de sa route courante
        (premier ou dernier élément)
  \item $j$ est une \textbf{extrémité} de sa route courante
\end{enumerate}

Si nécessaire, on retourne une route pour que les extrémités soient compatibles
(ex. : si $i$ est en début de route, on inverse pour le mettre en fin).

\subsection{Condition d'arrêt}

On s'arrête quand toutes les villes sont dans \textbf{une seule route},
qui forme alors le tour $0 \to v_1 \to v_2 \to \cdots \to v_{n-1} \to 0$.

\subsection{Complexité}

\begin{table}[h]
  \centering
  \begin{tabular}{lc}
    \toprule
    \textbf{Étape} & \textbf{Complexité} \\
    \midrule
    Calcul des savings       & $O(n^2)$ \\
    Tri des savings          & $O(n^2 \log n)$ \\
    Fusions (boucle greedy)  & $O(n^2)$ \\
    \midrule
    \textbf{Total}           & $O(n^2 \log n)$ \\
    \bottomrule
  \end{tabular}
  \caption{Complexité de l'algorithme Clarke-Wright}
\end{table}

% =============================================================================
\section{Algorithme complet}

\begin{algorithm}[H]
\caption{Clarke-Wright Savings Algorithm}
\begin{algorithmic}[1]
\State Choisir le dépôt $0$
\State Initialiser $n-1$ routes : $\{[i] : i \in V \setminus \{0\}\}$
\State Calculer $s(i,j) = c_{0i} + c_{0j} - c_{ij}$ pour tout $i \neq j$
\State Trier les savings par ordre décroissant
\For{chaque paire $(i,j)$ dans la liste triée}
  \If{$i$ et $j$ sont dans des routes différentes}
    \If{$i$ est extrémité de sa route \textbf{et} $j$ est extrémité de sa route}
      \State Orienter les routes pour que $i$ soit en fin et $j$ en début
      \State Fusionner : route$(i)$ + route$(j)$
    \EndIf
  \EndIf
\EndFor
\State Fermer le tour : relier la dernière ville au dépôt $0$
\end{algorithmic}
\end{algorithm}

% =============================================================================
\section{Explication intuitive}

Imagine que tu dois livrer $n$ colis depuis un entrepôt (dépôt).
Au départ, tu fais $n$ allers-retours séparés --- très coûteux.

Clarke-Wright te dit :
\begin{quote}
  \textit{"Regarde quelles paires de livraisons tu peux enchaîner directement
  sans repasser par l'entrepôt. Commence par celles qui te font économiser
  le plus de distance."}
\end{quote}

Le saving $s(i,j)$ mesure exactement ce gain : au lieu de faire
entrepôt $\to i \to$ entrepôt $\to j \to$ entrepôt, tu fais
entrepôt $\to i \to j \to$ entrepôt et tu économises $c_{0i} + c_{0j} - c_{ij}$.

\medskip
\textbf{Limites :} l'algorithme est glouton --- il prend la meilleure décision
locale à chaque étape sans reconsidérer les choix passés.
Il peut donc manquer l'optimum global.
En pratique, le gap avec l'optimal est de l'ordre de 10--15\%.

\medskip
\textbf{Comparaison avec les méthodes exactes :}

\begin{center}
\begin{tabular}{llll}
  \toprule
  & \textbf{ILP MTZ} & \textbf{Concorde} & \textbf{Clarke-Wright} \\
  \midrule
  Type           & Exact      & Exact       & Heuristique \\
  Optimalité     & Garantie   & Garantie    & Non garantie \\
  Complexité     & Exponentielle & Exponentielle & $O(n^2 \log n)$ \\
  Limite pratique & $n \leq 50$ & $n \leq 100\,000$ & $n > 1\,000\,000$ \\
  Gap typique    & 0\%        & 0\%         & 10--15\% \\
  \bottomrule
\end{tabular}
\end{center}

% =============================================================================
\section{Cas d'applications connus}

\begin{table}[h]
  \centering
  \begin{tabular}{ll}
    \toprule
    \textbf{Référence} & \textbf{Contribution} \\
    \midrule
    Clarke \& Wright (1964) & Article fondateur --- \textit{Scheduling of Vehicles from a Central Depot} \\
    VRP classique           & Baseline standard pour les problèmes de tournées de véhicules \\
    Logistique industrielle & Utilisé comme heuristique initiale dans de nombreux solveurs VRP \\
    \bottomrule
  \end{tabular}
  \caption{Références principales}
\end{table}

% =============================================================================
\section{Exemple de code}

\begin{lstlisting}[caption={Clarke-Wright Savings Algorithm --- Python}]
import math

def solve_tsp_clarke_wright(cities):
    """
    Clarke-Wright Savings Algorithm pour le TSP.
    cities : liste de tuples (x, y), cities[0] est le depot.
    Retourne : (tour, distance_totale)
    """
    n = len(cities)
    depot = 0

    def dist(i, j):
        dx = cities[i][0] - cities[j][0]
        dy = cities[i][1] - cities[j][1]
        return math.hypot(dx, dy)

    # Calcul des savings pour toutes les paires
    savings = []
    for i in range(1, n):
        for j in range(i + 1, n):
            s = dist(depot, i) + dist(depot, j) - dist(i, j)
            savings.append((s, i, j))
    savings.sort(reverse=True)

    # Initialisation : chaque ville est sa propre route [i]
    routes = {i: [i] for i in range(1, n)}
    # endpoint_to_route[v] = cle de la route dont v est une extremite
    endpoint_to_route = {i: i for i in range(1, n)}

    for s, i, j in savings:
        ri = endpoint_to_route.get(i)
        rj = endpoint_to_route.get(j)

        # i ou j n'est plus une extremite, ou deja dans la meme route
        if ri is None or rj is None or ri == rj:
            continue

        route_i = routes[ri]
        route_j = routes[rj]

        # Orienter route_i pour que i soit en fin
        if route_i[-1] == i:
            pass
        elif route_i[0] == i:
            route_i.reverse()
        else:
            continue  # i n'est pas extremite

        # Orienter route_j pour que j soit en debut
        if route_j[0] == j:
            pass
        elif route_j[-1] == j:
            route_j.reverse()
        else:
            continue  # j n'est pas extremite

        # Fusion : route_i + route_j
        new_route = route_i + route_j
        routes[ri] = new_route
        del routes[rj]

        # Mettre a jour les extremites
        head, tail = new_route[0], new_route[-1]
        # Supprimer les anciennes entrees pour i et j (plus extremites internes)
        endpoint_to_route.pop(i, None)
        endpoint_to_route.pop(j, None)
        endpoint_to_route[head] = ri
        endpoint_to_route[tail] = ri

    # Assembler le tour final : depot + route unique
    final_route = list(routes.values())[0] if routes else []
    tour = [depot] + final_route
    total = (sum(dist(tour[k], tour[k+1]) for k in range(len(tour)-1))
             + dist(tour[-1], tour[0]))
    return tour, total
\end{lstlisting}

% =============================================================================
\section{Benchmark TSP}

Mesuré sur graphes aléatoires (seed = 42, coordonnées uniformes dans $[0,100]^2$), CPU.
Le gap optimal est calculé par rapport à la solution ILP exacte (disponible seulement pour $n \leq 50$).
Modèle de complexité empirique : $t \approx 1{,}36 \times 10^{-6} \cdot n^{2{,}08}$.

\begin{table}[h]
  \centering
  \begin{tabular}{ccccc}
    \toprule
    \textbf{Nœuds} & \textbf{Distance} & \textbf{Temps CPU} & \textbf{Mémoire} & \textbf{Statut} \\
    \midrule
    1        & 0{,}00    & $< 0{,}001$\,s  & $<$\,1\,MB   & TRIVIAL  \\
    10       & 264{,}14  & $< 0{,}001$\,s  & $<$\,1\,MB   & OK       \\
    100      & 851{,}03  & 0{,}035\,s      & 0{,}4\,MB    & OK       \\
    1\,000   & 2\,512{,}54 & 1{,}93\,s     & 63\,MB       & OK       \\
    10\,000  & 8\,011{,}32 & 271\,s        & 6\,408\,MB   & LENT     \\
    100\,000 & ---       & ---             & ---          & SKIP (hors mémoire) \\
    \bottomrule
  \end{tabular}
  \caption{%
    Résultats benchmark Clarke-Wright (seed = 42).
    $n = 100\,000$ génère $\approx 5 \times 10^{9}$ paires de savings — hors mémoire avec l'implémentation $O(n^2)$.
  }
\end{table}

\medskip
\textbf{Prédictions (loi de puissance $t = 1{,}36 \times 10^{-6} \cdot n^{2{,}08}$) :}

\begin{center}
\begin{tabular}{cc}
  \toprule
  \textbf{$n$} & \textbf{Temps prédit} \\
  \midrule
  $1\,000\,000$  & $\approx 48\,days$ \\
  $10\,000\,000$ & $\approx 16\,years$ \\
  \bottomrule
\end{tabular}
\end{center}

\end{document}
